\chapter{投资组合与风险管理：从协方差估计到压力测试}
\label{ch:lower-portfolio-estimation}

\begin{itemize}
  % Generated from curriculum/manifest.json; do not edit by hand.
\item 直接先修：上册第 6 章《线性代数、分解、投影与条件数》、上册第 10 章《估计、检验、渐近、回归与多重检验》、上册第 14 章《浮点数、数值线代与病态问题》、上册第 16 章《研究有效性、样本外与交易摩擦》、上册第 17 章《可复现研究审计》。

  \item \textbf{学习产出：}构造并诊断风险模型，完成风险预算、Black--Litterman、CVaR 与稳健组合优化，再用尾部统计、非线性重估和反向压力寻找失败阈值。
  \item \textbf{研究边界：}风险矩阵、收益观点和尾部分布都是估计对象；优化器与风险数字不会修复错误数据、制度变化或不可成交权重。
\end{itemize}

收益共同波动决定组合风险。先把收益矩阵中心化，再看为什么协方差必半正定、何时必奇异，才能理解收缩为何改善某些方向的稳定性。


\section{先确定收益矩阵的含义}

令 $R\in\mathbb R^{T\times n}$ 的第 $t$ 行是同一可用时点的资产收益，列顺序在研究期内固定。样本协方差为
\[
 \widehat\Sigma=\frac{1}{T-1}\sum_{t=1}^{T}(r_t-\bar r)(r_t-\bar r)^\top.
\]
该公式看似简单，却依赖收益口径、时区、公司行动、缺失值和资产池协议。成对删除会让不同矩阵元素来自不同日期集合，拼出的矩阵甚至可能不再半正定。

% Generated by tools/build_figures.py; edit figures/figure-specs.json instead.
\MFQFigure{tex/figures/generated/lower-pf-risk-map.pdf}{协方差矩阵把权重方向映射为风险椭圆；主轴方向和长度分别对应特征向量与特征值。}{lower-pf-risk-map}


任意权重 $w$ 都满足
\[
 w^\top\Sigma w=\operatorname{Var}(w^\top r)\ge0,
\]
因此半正定不是数值偏好，而是协方差定义的必要条件。可逆也不等于可靠：若最小特征值很小，微小估计误差会被逆矩阵放大。报告至少应保存样本窗口、维度、最小特征值、条件数和缺失值策略。

\section{收缩与因子风险模型}

线性收缩写成
\[
 \widehat\Sigma_\lambda=(1-\lambda)\widehat\Sigma+\lambda T,
 \qquad 0\le\lambda\le1.
\]
目标 $T$ 可以是对角矩阵、等相关矩阵或结构模型。收缩用偏差交换方差，目的在于改善样本外稳定性，而不是让历史拟合更漂亮。教学实验以平均方差乘单位阵为目标，$\lambda=0.25$；透明公式与 sklearn \texttt{ShrunkCovariance} 完全一致，最小特征值为 $3.58\times10^{-5}$\cite{ledoit2004well}。

% Generated by tools/build_figures.py; edit figures/figure-specs.json instead.
\MFQFigure{tex/figures/generated/lower-pf-shrinkage.pdf}{收缩在高方差样本协方差与结构化目标之间权衡，通常以少量偏差换取更低估计波动。}{lower-pf-shrinkage}


因子模型把收益写成
\[
 r=Bf+\varepsilon,\qquad
 \Sigma=B\Sigma_fB^\top+D,
\]
其中需要 $\operatorname{Cov}(f,\varepsilon)=0$，且 $D$ 的对角元素非负。它减少需要估计的自由参数，也引入因子遗漏、暴露漂移与特异风险相关性等错设。事先确定三资产、两因子矩阵的迹为 $0.1013$；该值可从 $B,\Sigma_f,D$ 独立重建。

\begin{diagnostic}
协方差收缩和因子模型回答“如何约束估计误差”，并不回答“未来制度是否与估计窗口相同”。应把结构选择、强度选择和窗口选择写进样本外协议。
\end{diagnostic}

\section{边际风险、成分风险与风险平价}

组合方差 $v(w)=w^\top\Sigma w$ 的梯度为 $2\Sigma w$。采用方差口径时，第 $i$ 个资产的成分贡献定义为
\[
 \mathrm{RC}_i=w_i(\Sigma w)_i,
 \qquad \sum_i\mathrm{RC}_i=w^\top\Sigma w.
\]
若使用波动率口径，则还需除以 $\sqrt{w^\top\Sigma w}$；两种口径不能混写。等风险预算要求
\[
 \mathrm{RC}_1=\cdots=\mathrm{RC}_n.
\]

对角协方差的波动率为 $(0.2,0.3,0.4)$ 时，权重与逆波动率成正比，因此
\[
 w=\frac{(5,10/3,5/2)}{5+10/3+5/2}
   =\left(\frac6{13},\frac4{13},\frac3{13}\right).
\]
三个方差贡献恰好相等。仓库的透明实现求解严格凸问题
\[
 \min_{y_i>0}\left\{\frac12y^\top\Sigma y-\sum_i\log y_i\right\},
\]
其一阶条件 $y_i(\Sigma y)_i=1$ 在存在负相关时仍成立；Newton 法求出 $y$ 后再归一化。独立 SLSQP 直接最小化贡献差，最大权重差约 $5.8\times10^{-9}$。双实现一致只验证计算，不证明风险预算目标符合投资者负债或收益目标。

\MFQLead{估计误差必须进入报告}

若把 $\widehat\Sigma$ 当作真值，最终权重会显得过度精确。均值与协方差估计误差会被优化器放大，这是经典的研究有效性问题\cite{michaud1989markowitz}。教学实验对收益行做 500 次 IID bootstrap，每次重算组合波动率；点估计 $0.007950$，90\% 区间为 $[0.005984,0.009304]$。区间宽度提醒读者：优化器小数点后六位不等于输入有六位信息。

普通 IID bootstrap 需要独立同分布及统计量的相应正则性；只有可交换性不够。波动聚集、重叠收益、共同宏观冲击或横截面分组会破坏该假设，此时应改用区块、簇或状态分层 bootstrap。重采样单位必须与依赖结构一致。

协方差迹为 $7.415213\times10^{-5}$，权重为 $(0.434184,0.565816)$。这演示 provenance、时点协议和真实输入适配；正文明确禁止把两列宏观增长率冒充可交易资产收益或风险模型有效性证据。

本单元向本路线的综合项目交付风险矩阵、风险预算、估计区间、失败诊断和数据时点协议。

\MFQLead{样本协方差为何在高维下失效}

样本协方差
\[
S=\frac1{T-1}\sum_{t=1}^T(r_t-\bar r)(r_t-\bar r)^{\mathsf T}
\]
在 $N\ge T$ 时必然奇异；即使 $N<T$，最小特征值也可能主要是估计噪声。优化器会在这些看似低风险方向上建立巨大多空仓位。报告维数比 $N/T$、条件数和特征值谱比只报“矩阵可逆”更有信息。

收缩估计 $\hat\Sigma=(1-\delta)S+\delta F$ 把样本矩阵拉向结构化目标 $F$。目标可取常相关、对角或因子模型。$\delta$ 越大方差越低、结构偏差越高；应通过解析收缩或训练期验证选择，而不是为得到漂亮权重手调。

\MFQLead{因子与特异风险的归因}
组合因子暴露为 $B^Tw$，在因子与特异项不相关的假设下，总方差分解为 $(B^Tw)^T\Sigma_f(B^Tw)+w^TDw$。若特异项相关，$D$ 应包含这些非对角协方差；把它强行设成对角会改变风险模型。

\MFQLead{动态协方差与状态依赖}

固定窗口协方差假设窗口内稳定。指数加权、DCC-GARCH 或状态条件协方差允许时变，但引入更多参数和模型风险。危机期相关通常上升，使用平静期矩阵会低估压力风险。

在决策时只能使用此前收益更新矩阵。若当天收盘组合使用包含当天收盘收益的协方差并按同一收盘价交易，需明确是否有足够计算和成交时间，否则应滞后一日。

这些估计对象将在组合风险路线 Capstone 中与优化、尾部压力和限制报告放在同一条证据链上。
\subsection{秩、收缩与风险贡献的同一例子}

记中心化矩阵为 $X$，则 $S=X^TX/(T-1)$。任意 $w$ 有 $w^TSw=\|Xw\|^2/(T-1)\geq0$；且中心化使行和为零，所以 $\operatorname{rank}S\leq\min(N,T-1)$。这是 $N\geq T$ 时必奇异的直接证明。

取 $S=\left(\begin{smallmatrix}1&1\\1&1\end{smallmatrix}\right)$，特征值为 2 和 0。向单位阵收缩得到 $S_\lambda=(1-\lambda)S+\lambda I$，特征值为 $2-\lambda,\lambda$；$0<\lambda\leq1$ 时可逆，条件数为 $(2-\lambda)/\lambda$。但低方差方向被抬高来自结构假设，不是新增观测。

对 $\Sigma=\operatorname{diag}(0.04,0.16)$，逆波动权重为 $(2/3,1/3)$，方差贡献分别为 $0.04(2/3)^2=0.16(1/3)^2=4/225$。总波动为 $\sqrt{8/225}$，波动贡献各为其一半。notebook 同时画出收缩后的谱与权重，并通过重新抽样观察估计变化。
